\scanpage{89}\mspage{33}
% L1-3
des deux catégories donnent naissance \inserted{chacun} à \deleted{\unsure{des}} \unsure{foncteurs|foncteur} \unsure{contravariants|contravariant} sur $S$
\[
 \bar{\xi}:S^{\unsure{\mathrm o|\circ}}\longrightarrow\mathrm{Hot}_W,
\]
% L4-6
\inserted{ceux de} la seconde (qui n'est autre que $(\mathrm{Fibsc}(S^{\mathrm o}))^{\mathrm o}$) donnent naissance \inserted{chacun} à \deleted{\unsure{des}} foncteurs \uline{covariants}
\[
 \bar{\eta}:S\longrightarrow\mathrm{Hot}_W,
\]
% L7-11
ce qui n'est pas pareil ! Ce n'est ``pareil'' que si ce sont des coeff. locaux, ce qui signifie exactement que $\xi$ est dans $\underline{\mathrm{Fibsc}}_0\,S$ resp. $\eta$ dans $\deleted{\unsure{\mathrm{Fibsc}}}\underline{\mathrm{Cofibsc}}_0\,\deleted{\unsure{\xi}}S$. On s'attend donc bien \textcircled{\scriptsize 1\textsuperscript{o}} que les flèches
% L12
\[
 \underline{\mathrm{Fibsc}}_0\,\deleted{\unsure{(S)}}\unsure{X|S}\longrightarrow\underline{\mathrm{Fib}}_0\,S,\quad
 \underline{\mathrm{Cofibsc}}_0\,\deleted{\unsure{(}}S\deleted{\unsure{)}}\longrightarrow\underline{\mathrm{Cofib}}_0\,S
\]
% L13-15
induisent des équivalences sur les localisées (chose que j'ai été incapable de prouver tantôt)$^{*}$. \margin{$^{*}$ Les \unsure{quasi-inv.} étant bel et bien donnés par $\Phi_0,\Psi_0$ -- car par quoi d'autre ?}
% slanted
Par contre, il est \inserted{pire que} déraisonnable de s'attendre \inserted{ici} qu'il en soit de $=$ pour les foncteurs
% L16-17
\[
 \underline{\mathrm{Fibsc}}_0\,S\longrightarrow\underline{\mathrm{Fibsc}}\,S,\quad
 \underline{\mathrm{Cofibsc}}_0\,S\unsure{\mathord{-}|\longrightarrow}\underline{\mathrm{Cofibsc}}\,S
\]
% L18-21
(ça devrait donner par localisation des foncteurs pleinement fidèles d'inclusion, des catégories des objets loc. const. dans $\mathrm{HOT}_W(S)$ resp. dans $(\unsure{\mathrm{HOt}|\mathrm{HOT}}_W(S^{\mathrm o}))^{\mathrm o}$. Et de même pour
% archive sheet number: 88
